\section{Introduction}\label{sec:introduction}

Euler's totient function can take the same value at many different
integers. Write
\[
 g(n)=\#\{m\ge1:\varphi(m)=n\}
\]
for the size of its fiber at $n$. We prove the following result.

\begin{theorem}\label{thm:totient}
For every real $\eps>0$, there are infinitely many positive integers
$n$ such that
\[
 g(n)>n^{1-\eps}.
\]
\end{theorem}

\Cref{thm:totient} resolves positively Erd\H{o}s's conjecture on the
largest fibers of Euler's totient function. Pomerance
\cite[p.~84]{Pomerance1980} formulates the conjecture as $C=1$, where
$C$ is the supremum of the exponents $c$ for which $g(n)>n^c$ infinitely
often, and attributes it to Erd\H{o}s \cite{Erdos1956}. An elementary
upper bound is $g(n)\ll_\eta n^{1+\eta}$ for every $\eta>0$; we recall
it in \Cref{sec:totients}. Thus the power exponent of $n$ in the theorem is
optimal.

The arithmetic input is a result about the prime factors of the
predecessor of a prime. For $m>1$, let $P^+(m)$ be its largest prime
factor. Erd\H{o}s's work on totient multiplicities and smooth shifted
primes dates to \cite[Section~3]{Erdos1935}. The underlying
construction takes products of many primes whose predecessors have
few available prime factors: many distinct products then have the same
totient. The expectation that $p-1$ can have
all its prime factors smaller than every fixed power of $p$ was
recorded in \cite{Erdos1956}; see also
\cite[Introduction]{Lichtman2022}. We prove a quantitative form.

\begin{theorem}\label{thm:smooth}
For every fixed $0<\delta<1/4$, as $x\to\infty$,
\begin{equation}\label{eq:main-smooth}
 \#\{p:\ p\text{ prime},\ 2x<p\le5x,
                   \ P^+(p-1)\le x^\delta\}
 \ge x^{1-o(1)}.
\end{equation}
The quantity $o(1)$ may depend on $\delta$.
\end{theorem}

In particular, for every $\eps>0$ there are infinitely many primes
$p$ with $P^+(p-1)\le p^\eps$. This resolves the smooth-predecessor
conjecture positively as well. The count in \Cref{eq:main-smooth}
also holds for every fixed $\delta>0$, by using a smaller positive
exponent when necessary. No uniformity as $\delta\to0$ is
asserted in \Cref{thm:smooth} or needed for \Cref{thm:totient}.

There is a substantial literature on fixed smoothness exponents.
Baker and Harman \cite[Section~1, Theorem~1 and Corollary~1]{BakerHarman1998}
obtained exponent $0.2961$ for shifted primes, with the corresponding
totient multiplicity exponent $0.7039$. This improved Friedlander's
threshold $1/(2\sqrt e)+\eps$; their introduction also describes
preceding work of Pomerance, Balog, and Fouvry--Grupp.
Lichtman \cite[Theorem~1.1 and Corollary~1.3]{Lichtman2022}
proved a lower bound $x/(\log x)^C$ for $x<p\le2x$ and
$P^+(p-1)\le x^\beta$ whenever
\[
 \beta>\frac{15}{32\sqrt e}=0.284311\ldots,
\]
and obtained multiplicity exponent $0.7156$.
These results concern fixed positive exponents. The expected
Dickman law predicts a positive proportion of primes for every such
exponent; see Granville \cite[Section~5.3]{Granville2008}.
Bharadwaj and Rodgers \cite[Theorem~7, Proposition~2,
and Conjecture~5]{BharadwajRodgers2026} prove the full
Poisson--Dirichlet law for sequences satisfying their regularity and
congruence-uniformity conditions together with level-one distribution.
For shifted primes this gives the prediction under the
Elliott--Halberstam conjecture; their Proposition~2 establishes
distribution level one half and the other conditions unconditionally.
Our unconditional bound $x^{1-o(1)}$ gives the full power exponent of
the count for every fixed positive smoothness exponent. A positive
limiting proportion is a stronger conclusion.
The final passage from plentiful smooth shifted primes to large
totient fibers belongs to the method of Erd\H{o}s and Pomerance;
Pomerance states the relevant transfer explicitly in
\cite[Theorem~B]{Pomerance1980}. We give the needed
product-and-pigeonhole proof in full.

Smooth predecessors also enter the construction of Carmichael numbers
by Alford, Granville and Pomerance \cite{AlfordGranvillePomerance1994},
where their supply is combined with a separate distribution theorem
for primes in arithmetic progressions. This is one reason that
estimates for shifted smoothness have applications beyond the
factorization problem itself. The contribution here includes the
general graph and kernel theorems that produce the arithmetic lower
bound, as well as the optimal exponent for totient fibers.

These graph and ideal-kernel estimates also serve as inputs to the
companion paper \cite[Theorem~1.1]{PrimePredecessors}, where a separate
determinant-graph estimate and prime-extraction argument yield the full
Poisson--Dirichlet law for prime predecessors.

\subsection{The analytic mechanism}

The main obstacle is to detect primality in a sequence whose
predecessors have a strongly constrained factorization. Congruence
counts provide an initial sieve, and a further bilinear cancellation
estimate permits the composite contribution to be removed. This
division of work is familiar from prime-detecting sieves; compare
Friedlander and Iwaniec \cite[Section~1]{FriedlanderIwaniec1998}.
Our coefficient and range hypotheses are stated and verified below
for the particular weight used here.

We count primes of the form $2u+1$, where $u$ has a prescribed
factorization over many logarithmic scales. Some factors lie in
two separated ranges of subpower-sized primes, divided into small and
big groups. Their weights mark distinct prime divisors. The functions
attached to an integer depend only on the part of $u$ outside these
groups; the marked weights supply the divisibility
labels for the dilation graph. Other factors
occupy a geometric sequence of size bands, which permits a divisor
of $u$ to be chosen close to a specified size. One band contains a
long factor ranging over rough integers, meaning integers with no
prime factor below a prescribed cutoff.

The main new analytic ingredient is a transference theorem for
weighted dilation graphs. A physical state consists of an integer and
ordered lists of distinct group-prime divisors. An edge has shift
$kD$ for an integer $k\ne0$, where $D$ is a product of primes shared
by its endpoint lists,
together with independent free prime factors. The corresponding ideal
operator acts on the big-group lists: its labels are independent, with
probabilities proportional to $1/p$, and repetitions are allowed. The list length is sufficiently large
but fixed as $x$ tends to infinity. A small ideal operator norm yields
an averaged high moment of the physical graph, and hence cancellation
in pairings whose first endpoint is independent of its mark list.

The transference proof averages the integer root by the Chinese
remainder theorem. Repeated divisibility queries then pay a reciprocal
prime factor only once; a memory retains their congruences between
active uses. Averaging the omitted small marks makes transfers between
memory and the active lists rare. A rank expansion restores distinct
prime births while preserving cancellation on most edges: many
independent birth equalities supply reciprocal-prime savings, whereas
a small equality rank affects only a small fraction of the edges.
The theorem applies to arbitrary signed kernels satisfying its explicit
complexity and norm hypotheses.

Related divisibility-graph methods have a substantial history.
Matom\"aki, Radziwi\l\l{} and Tao
\cite[Theorem~5.1]{MatomakiRadziwillTao2016} studied connectivity
of a graph with prime-divisibility edges. Tao's entropy-decrement
argument \cite{Tao2016} replaces such divisibility conditions by
their mean at a suitable scale. Helfgott and Radziwi\l\l{}
\cite[Section~1.5]{HelfgottRadziwill2021} study centered adjacency
operators through signed closed-walk moments, with repeated primes
creating dependent congruences. Pilatte
\cite[Theorem~2.4 and Section~5]{Pilatte2026} develops this approach
for shifts that are products of primes. These works explain the
roles of centering, long moments, and repeated-label bookkeeping.
Our graph also carries ordered active lists, and its comparison
with an independent-label operator requires the explicit lifespan
and factorial-memory construction proved in \Cref{sec:transference}.

The independent-label norm is made small by a comparison construction.
Logarithmic localization and character kernels allow shared prime
products to be replaced by independent products. The big groups are
split into two blocks, with a bounded number of selected coordinates
in each group. After averaging the list coordinates unused by the
multiplier, a coordinate decomposition isolates components that have
mean zero in all selected coordinates of one block. Permutation
symmetry makes these components small, and signed comparison terms
cancel the other components up to a controlled error. For each
prescribed fixed accuracy, the construction gives kernels uniform in
the additive frequency, with complexity fixed
before the length of the marked lists is chosen.

Applying transference to the multiplicatively invariant endpoints
reduces a fixed shifted correlation to averages over larger shifts.
Products of free primes control the minor arcs. On the major arcs,
a marked small prime and the compulsory long rough-integer factor
control the local Fourier energy. A low-complexity character and
Mellin discrepancy condition supplies the remaining cancellation.
The argument permits arbitrary bounded residual coefficients.

These shifted correlations yield a Type II estimate for
$mn=2u+1$. A probabilistic split of the geometric prime bands chooses
a divisor $e$ of $u$ close to the scale of $m$. After Cauchy--Schwarz,
off-diagonal factorizations become shifted endpoints by an exact
determinant identity. The Type II estimate then permits primality
and roughness tests on cofactors to be replaced by elementary rough
proxies. Congruence estimates and a weighted sieve remove the
composite values of $2u+1$; a separate upper bound treats the small
balanced range.

\subsection{Organization and dependencies}

\Cref{sec:preliminaries} fixes the conventions and proves the
elementary sieve used later. The central transference estimate is
proved in \Cref{sec:transference}, and the independent-label kernels
are constructed in \Cref{sec:ideal}. Together they lead to the
shifted-correlation theorem in \Cref{sec:shifts}. The Type II
reduction occupies \Cref{sec:typeii}. The prime extraction and all
its parameter choices are completed in \Cref{sec:primes}.
\Cref{sec:totients} finishes the proof of \Cref{thm:totient}.
The main implications are summarized in \Cref{fig:proof-dependencies}.

The sieve geometry is fixed first, the required analytic accuracies
next, and the discrepancy precision of the cofactor proxies last.
The relevant theorems state this choice order, and the prime extraction
verifies it. In particular, the needed discrepancy is established
independently of the estimate that uses it.

\begin{figure}[!ht]
\centering
\begin{tikzpicture}[>=Stealth,
  box/.style={draw=blue!45!black,fill=blue!3,rounded corners=2pt,
              align=center,text width=3.15cm,minimum height=.78cm,
              inner sep=4pt,font=\small},
  link/.style={->,thick,draw=blue!55!black}]
 \node[box] (transfer) at (0,0)
   {Transference\\\Cref{thm:transference}};
 \node[box] (ideal) at (3.8,0)
   {Small ideal norms\\\Cref{thm:ideal}};
 \node[box] (shift) at (1.9,-1.65)
   {Shift cancellation\\\Cref{thm:shift}};
 \node[box] (typeii) at (6,-1.65)
   {Type II estimate\\\Cref{thm:typeii}};
 \node[box] (sieve) at (7.6,0)
   {Candidate, Type I,\\and sieve bounds\\\Cref{sec:primes}};
 \node[box] (smooth) at (11.4,-1.1)
   {Smooth shifted primes\\\Cref{thm:smooth}};
 \node[box] (totient) at (11.4,-2.6)
   {Large totient fibers\\\Cref{thm:totient}};
 \draw[link] (transfer)--(shift);
 \draw[link] (ideal)--(shift);
 \draw[link] (shift)--(typeii);
 \draw[link] (typeii)--(smooth);
 \draw[link] (sieve)--(smooth);
 \draw[link] (smooth)--(totient);
\end{tikzpicture}
\caption{Main result dependencies. The auxiliary branch includes
candidate mass, Type I distribution, proxy discrepancy, the elementary
sieve, and the separate balanced-range bound. Arrows record the implications used
in the proof; the order of choosing constants is described in the text.}
\label{fig:proof-dependencies}
\end{figure}
